\section{Positive algebraic dimension}
\label{pos:section}

\begin{proposition}\label{pos:nonvanishing}
Assume Assumptions~\ref{hyp:campana}, \ref{hyp:mmp}, and
\ref{hyp:effective-abundance}. Let \(W\) be a smooth non-uniruled
compact Kähler fourfold with \(q(W)=0\) and \(0<a(W)<4\). Then
\(\kappa(W,K_W)\geq0\).
\end{proposition}

The first step constructs an actual canonical pullback over a
projective base. We then prove nonvanishing of that base line in each
of its three possible dimensions.

\subsection{A contraction index bound}

\begin{lemma}[A local rationality bound]\label{pos:rationality}
Let \((Z,\Delta)\) be an ordinary rational klt pair on a normal
\(n\)-dimensional compact Kähler space, with actual \(\Q\)-Cartier
adjoint \(J=K_Z+\Delta\). Let \(\pi:Z\to T\) be a projective
bimeromorphic morphism with connected fibers to a normal compact
Kähler space. Assume that \(-J\) is relatively ample and that all
contracted curves have class in one \(J\)-negative ray \(R\) of
\(\NA(Z)\). Let \(k>0\) be an integer such that \(kJ\) is Cartier,
and let \(L\) be a line bundle with \(L\cdot R>0\). For any
nontrivial fiber \(F\) of \(\pi\),
\begin{equation}\label{pos:index-bound}
 r_F:=\sup\{u:L+uJ\text{ has nonnegative degree on all curves of }F\}
 \geq\frac{1}{k(n+1)}.
\end{equation}
No local analytic \(\Q\)-factoriality is required.
In particular, this applies to the contractions in
Assumption~\ref{hyp:mmp}, with \(n=4\).
\end{lemma}

\begin{proof}
Every curve \(C\) in \(F\) has nonzero class, since a Kähler class
has positive degree on it, and its class belongs to \(R\). Therefore
\[
 r_F=\frac{L\cdot C}{-J\cdot C}>0
\]
is independent of \(C\). The restriction of \(L\) to the projective
fiber is numerically a positive multiple of the relatively ample
rational line \(-J\). It is thus ample on \(F\), and the openness of
fiberwise ampleness makes \(L\) relatively ample after shrinking around
the image point of \(F\).

Work over a small Stein neighborhood of that point, with compactum
equal to the point. At \(r_F\), the restriction of \(L+r_FJ\) is
numerically trivial, hence nef but not ample on the nontrivial fiber.
Fujino's relative rationality theorem applies to this finite positive
threshold: in lowest terms its denominator is at most
\(k(\dim F+1)\) \cite[Theorem~4.3.1]{FujinoAnalyticCone}. The
non-lc-locus condition is empty for a klt pair, and the singleton
compactum satisfies the required local condition on the Stein base.
The relative numerical space over it is finite dimensional, as is also
seen by restriction to the projective fiber. The theorem consequently
gives \(r_F\geq1/(k(\dim F+1))\), which implies
\eqref{pos:index-bound}.

The line-bundle formulation is expressly permitted by
\cite[Remark~4.3.4]{FujinoAnalyticCone}; the nearby ampleness assertion
is \cite[Lemma~2.2.4]{FujinoAnalyticCone}. Canonical divisors in that
theorem may be handled locally, in its formal canonical-class
convention. Alternatively, over the Stein neighborhood the proper
direct image of a rank-one canonical sheaf is coherent of rank one,
since the contraction is bimeromorphic. Cartan's theorem supplies a
generically nonzero section, hence a meromorphic canonical
representative there. This requires no global meromorphic canonical
frame on the compact source and no local \(\Q\)-factoriality. The
Cartier index used in the bound belongs to this stage alone.
\end{proof}

\subsection{An actual canonical pullback}

Put \(d=a(W)\). Resolve a map defined by \(d\) algebraically independent
meromorphic functions on \(W\). Stein factorization has a projective
base, since its finite map has projective image. Resolving that base
and the main graph, we may replace \(W\) by a smooth compact Kähler
model on which there is a fibration
\[
 b:W\longrightarrow S_0,
 \qquad S_0\text{ smooth connected projective},\quad\dim S_0=d.
\]
The generic fibers remain connected through these modifications; Stein
factorization over the normal resolved base gives connected fibers
everywhere. Fix a very ample line bundle \(H\) on \(S_0\).

Since \(K_W\) is pseudo-effective and \(4-d\leq3\), the restriction
and parameter argument in Section~\ref{red:section} gives
\(\rk b_*\OO_W(\ell K_W)>0\) for some positive integer \(\ell\).
Coherence and ample twisting on the projective base imply that
\(b_*\OO_W(\ell K_W)\otimes H^{\ell N}\) has a nonzero global
section for every sufficiently large integer \(N\). Thus there is
\(N_0\) such that
\begin{equation}\label{pos:all-large-twists}
 \kappa(K_W+Nb^*H)\geq0
 \qquad\text{for every integer }N\geq N_0.
\end{equation}

\begin{proposition}[The pullback model]\label{pos:pullback-model}
There are a canonical globally strongly \(\Q\)-factorial compact
Kähler fourfold \(V\), a normal projective variety \(T\) of dimension
\(d\), a fibration \(g:V\to T\), and an actual rational line bundle
\(A_T\in\Pic(T)\otimes\Q\), such that
\begin{equation}\label{pos:canonical-pullback}
 K_V\sim_{\Q}g^*A_T.
\end{equation}
The space \(V\) is obtained from the current smooth \(W\) by a finite
empty-boundary \(K\)-program allowed by Assumption~\ref{hyp:mmp}, and
\(K_V\) remains pseudo-effective.
\end{proposition}

\begin{proof}
At a stage \(V_i\) of the empty-boundary program that still maps to
\(S_0\), write \(b_i:V_i\to S_0\) and \(H_i=b_i^*H\).
All these stages are canonical. Indeed, a common projective resolution
of a negative canonical step gives the comparison
\begin{equation}\label{pos:canonical-step-comparison}
 p^*K_{\mathrm{before}}\sim_{\Q}q^*K_{\mathrm{after}}+G,
 \qquad G\geq0,
\end{equation}
with \(G\) exceptional over the after model. The comparison follows
from the projective analytic negativity lemma over the contraction
base: there is no extraction, and canonical negativity supplies the
relative sign. For flips both sides and their common resolution are
projective over that base. Compatible local canonical comparisons
glue to the intrinsic discrepancy divisor in
\eqref{pos:canonical-step-comparison}. Discrepancies therefore do not
decrease. Starting from smooth \(W\), this preserves canonicity. The
remaining category and pseudo-effectivity properties are part of
Assumption~\ref{hyp:mmp}.

Choose a Cartier index \(k_i\) for \(K_{V_i}\), and choose an integer
\begin{equation}\label{pos:stage-choice}
 N\geq N_0,\qquad N>5k_i.
\end{equation}
This choice is made anew at the current stage. The actual line
\(K_{V_i}+NH_i\) has an effective rational klt boundary representative:
take a general divisor in a sufficiently high free multiple of
\(NH_i\) and divide by that multiple. On a fixed log resolution the
system has no fixed exceptional component, and Bertini makes the
chosen member transverse to the exceptional strata; its small
coefficient preserves klt. Its adjoint is pseudo-effective because
\(K_{V_i}\) is pseudo-effective and \(H_i\) is semipositive.

If this adjoint is not analytically nef, Assumption~\ref{hyp:mmp}
supplies an extremal ray on which \(K_{V_i}+NH_i\) is negative.
Since \(H_i\) is the pullback of a semipositive form, the ray is also
\(K_{V_i}\)-negative. Take its contraction and its negative step for
the empty-boundary program. In fact
\begin{equation}\label{pos:base-line-zero-on-ray}
 H_i\cdot R_i=0.
\end{equation}
Otherwise Lemma~\ref{pos:rationality}, with \(L=H_i\) and
\(J=K_{V_i}\), gives
\[
 \frac{1}{5k_i}
 \leq\frac{H_i\cdot C}{-K_{V_i}\cdot C}
 <\frac1N
\]
on any curve of a nontrivial contraction fiber, contrary to
\eqref{pos:stage-choice}.

Every connected projective fiber of the contraction maps to a point
under \(b_i\). If its image had positive dimension, some curve in that
fiber would map nontrivially to the projective base and have positive
\(H_i\)-degree. Proper descent over the normal contraction target
therefore factors \(b_i\) through it. In a flip, composition with the
positive-side morphism gives \(b_{i+1}\). Thus the actual pullback line
\(H_i\), not merely its numerical class, continues to be the pullback
of the fixed \(H\).

Repeat this procedure as long as the selected twist is not nef.
Every step actually taken is a step of the single empty-boundary
\(K\)-program starting from \(W\). The changing integers \(N\)
select rays; they do not change that program's boundary. If the
procedure were infinite, it would contradict the arbitrary termination
assertion in Assumption~\ref{hyp:mmp}. It therefore reaches a stage
\(V_i\) at which its selected \(K_{V_i}+NH_i\) is nef. In particular,
if the canonical program reaches a nef canonical class, the selected
semipositive twist is nef there as well.

For this final value of \(N\), \eqref{pos:all-large-twists} already
gave a section on the original \(W\). Clear the finitely many Cartier
indices of the chosen program. Since no step extracts a prime and all
maps agree to \(S_0\), that section pushes through every step as a
section of the corresponding actual twisted canonical power.
Reflexive extension across codimension two on each normal target
justifies the pushforward. Hence
\(\kappa(K_{V_i}+NH_i)\geq0\) before abundance is invoked.
Assumption~\ref{hyp:effective-abundance}, applied to the klt
representative above, makes this actual line semiample.

Set \(V=V_i\), choose a free integral multiple
\(m(K_V+NH_i)\), and let \(h:V\to\PP^r\) be its morphism.
Take the Stein factorization \(g:V\to T\) of \((b_i,h)\).
The variety \(T\) is normal and finite over the projective image,
hence projective. Let \(a:T\to S_0\) and \(c:T\to\PP^r\) be the
induced maps. We have \(\dim T\geq d\) because \(a\) is surjective,
and \(\dim T\leq a(V)=d\) because \(T\) is projective. Define
\begin{equation}\label{pos:actual-base-line}
 A_T=\frac1m c^*\OO_{\PP^r}(1)-N a^*H
 \quad\text{in }\Pic(T)\otimes\Q.
\end{equation}
The defining evaluation isomorphism for \(h\) and the actual identity
\(H_i=g^*a^*H\) prove \eqref{pos:canonical-pullback}.
\end{proof}

Choose a smooth projective resolution \(\mu:S\to T\), and resolve
the main graph to obtain a smooth compact Kähler \(M\) with
\(p:M\to V\) and a fibration \(f:M\to S\). Put \(A=\mu^*A_T\).
Canonicity and the fixed actual isomorphism in
\eqref{pos:canonical-pullback} give
\begin{equation}\label{pos:resolved-pullback}
 K_M\sim_{\Q}f^*A+R,
 \qquad R\geq0\text{ and }R\text{ is }p\text{-exceptional}.
\end{equation}
This is an identity of rational line bundles; all further resolutions
use its canonical transform. Pullback of holomorphic one-forms gives
\(q(S)\leq q(M)=q(W)=0\).

The rational line \(A\) is pseudo-effective. To see this with the
analytic definition, pull a positive representative \(\Theta\) of
\(c_1(K_V)\), with local potentials, to \(M\). The map \(p\) is
dominant, so its plurisubharmonic potentials do not become identically
minus infinity. If \(\omega\) is a Kähler form on \(M\), then
\[
 f_*(p^*\Theta\wedge\omega^{4-d})
\]
is a closed positive \((1,1)\)-current. The projection formula puts its
class in \(v\,c_1(A)\), where the degree-zero closed current
\(f_*(\omega^{4-d})=v\) is the positive constant volume of a smooth
general fiber. Division by \(v\) proves the assertion. Since \(S\)
is projective, this also gives numerical divisor pseudo-effectivity
\cite{BDPP}.

It remains to prove \(\kappa(S,A)\geq0\). A section then pulls back
through \eqref{pos:resolved-pullback}, multiplied by the effective
exceptional section, to a canonical plurisection on \(M\), and hence
on \(W\). If \(d=1\), the smooth projective curve \(S\) has
\(q(S)=0\), so it is \(\PP^1\); a pseudo-effective rational line on
it has nonnegative degree and has a section at a divisible multiple.
The two remaining dimensions require more information from the
fibration.

\subsection{A surface base and fiber powers}

Assume \(d=2\). Fix a sufficiently divisible positive integer \(m_0\)
for \eqref{pos:resolved-pullback}. On very general smooth fibers one
has
\begin{equation}\label{pos:surface-fiber-rank}
 h^0(F,mK_F)=1\qquad(m>0,\ m_0\mid m).
\end{equation}
At least one section is supplied by the effective divisor \(R|_F\).
If for a fixed divisible \(m\) the generic dimension were at least
two, generic base change and ample twisting of
\(f_*\OO_M(mK_M)\) would give two sections whose ratio is nonconstant
on a general fiber. Together with meromorphic functions from the
two-dimensional projective base this would give algebraic dimension
at least three, contrary to \(a(M)=2\). The analytic jumping loci
for the countably many \(m\)'s can be excluded simultaneously, proving
\eqref{pos:surface-fiber-rank}.

\begin{lemma}[A base intersection inequality]\label{pos:surface-inequality}
There are nonnegative rational numbers \(c_P\), indexed by the
finitely many \(\mu\)-exceptional curves on \(S\), such that
\begin{equation}\label{pos:surface-intersection}
 \left(A-K_S+\sum_Pc_PP\right)\cdot H'\geq0
\end{equation}
for every ample divisor \(H'\) on \(S\).
\end{lemma}

\begin{proof}
Choose a dense open set of \(S\) over which \(f\) is smooth and the
generator supplied by \eqref{pos:resolved-pullback} is not identically
zero on any fiber. This requires deleting a proper analytic subset:
a horizontal divisor cannot contain a whole fiber over a divisor in
the base without being vertical, by the fiber-dimension theorem; the
remaining degeneracies are proper analytic images or jumping loci.
Include all \(\mu\)-exceptional curves among the finitely many
complementary curves.

For each such curve \(P\), fix a component \(D\) of \(f^*P\)
dominating \(P\), write
\[
 e=\ord_D(f^*P),\qquad h=\operatorname{coeff}_D R,
\]
and make these choices before choosing any ample test curve or fiber
power. If \(P\) is not exceptional over \(T\), take \(D\) with
\(h=0\). Indeed, over a general point of its image divisor in the
normal \(T\), a local equation pulls back under the holomorphic
\(g\) to a divisor on \(V\) dominating that base divisor. The strict
transform of one such component has zero coefficient in the
\(p\)-exceptional \(R\). For an exceptional \(P\), set
\begin{equation}\label{pos:exceptional-pole-allowance}
 c_P=h/e.
\end{equation}

Replace \(H'\) by a sufficiently large very ample multiple and choose
a smooth member \(C\) of genus at least one. It can be chosen
transverse to all complementary curves at their general points,
avoiding their finitely many excluded points and any isolated bad
base points. Moreover, \(M_C=f^{-1}(C)\) is smooth by Bertini for
the pulled-back free system, and connected because \(f\) has connected
fibers. To retain \eqref{pos:surface-fiber-rank}, choose \(C\) through
one parameter where all those countably many equalities hold. A high
enough linear system through that point is free away from it; on its
fiber \(f\) is a submersion, so the same Bertini conclusion holds
there. Each jumping locus then cuts a proper analytic subset of
\(C\). Thus very general fibers of \(f_C:M_C\to C\) satisfy all
the equalities in \eqref{pos:surface-fiber-rank}.

By adjunction,
\begin{equation}\label{pos:curve-relative-canonical}
 K_{M_C/C}\sim_{\Q} f_C^*((A-K_S)|_C)+R_C,
 \qquad R_C=R|_{M_C}\geq0.
\end{equation}
At each \(z\in C\cap P\), a general point of the chosen component
\(D\) over \(z\) admits coordinates \((x,y_1,y_2)\) on \(M_C\)
and a parameter \(t\) on \(C\) in which
\begin{equation}\label{pos:surface-local-chart}
 t=x^e,\qquad R_C=h\{x=0\}
\end{equation}
in that chart. A unit in the pullback equation is absorbed into \(x\).
Such charts exist over general points of \(P\): the map \(D\to P\)
is generically submersive, and its intersections with the other
divisorial supports have smaller generic fiber dimension. Properness
and the fiber-dimension theorem allow the exceptional base points to
be discarded before \(C\) is chosen. These choices are independent of
the next integer \(r\).

For \(r\geq1\), take the main component of the \(r\)-fold fiber
product of \(M_C\) over \(C\), and a smooth compact Kähler resolution
\(Z_r\) of it. Over the good open set this fiber product is smooth
with connected fibers \(F^r\), hence has a unique main component.
It is an analytic subspace of a product of compact Kähler spaces;
the required Kähler resolution and connected-fiber morphism to \(C\)
follow as above. The general fiber has a canonical section by
\eqref{pos:surface-fiber-rank}, and \(\kappa(C)\geq0\).
Assumption~\ref{hyp:campana}, with zero boundary and in dimension
\(2r+1\), gives a nonzero section of \(mK_{Z_r}\) for some \(m>0\).
By taking a power, make \(m\) divisible by \(m_0\) and all indices
in use. It may depend on \(r\) and \(C\).

Over good parameters, the product formula and
\eqref{pos:surface-fiber-rank} make the fiberwise pluriform space
one-dimensional. The chosen section is therefore the product of the
\(r\) relative generators in
\eqref{pos:curve-relative-canonical}, times a base pluriform and a
scalar base coefficient. This coefficient is a section on the good
open of the line
\begin{equation}\label{pos:fiber-power-coefficient-line}
 \OO_C\bigl(m(K_C+r(A-K_S)|_C)\bigr).
\end{equation}
Indeed, its quotient by the displayed product is constant on very
general fibers; local submersion sections and analytic continuation
give meromorphic descent to that open. Evaluation at points where the
product does not vanish identically on the fiber makes the descended
coefficient holomorphic there.

We compute its order at a missing point \(z\). In
\eqref{pos:surface-local-chart}, a frame of
\(\OO_C(m(A-K_S)|_C)\) maps, up to a unit, to
\[
 x^{mh}\left(\frac{dx\wedge dy_1\wedge dy_2}{dt}\right)^{\!m}.
\]
On a normalization branch of the local \(r\)-fold product we have
\(x_1=x\), \(x_j=\zeta_jx\), where \(\zeta_j^e=1\). Multiplying
the \(r\) relative factors by one base factor \((dt)^m\), and using
\(dt=ex^{e-1}dx\), gives order
\begin{equation}\label{pos:fiber-power-order}
 m\bigl(rh-(r-1)(e-1)\bigr)
\end{equation}
along \(x=0\). This smooth normalization branch belongs to the main
component, being a closure of points with \(t\ne0\). The global
pluriform on \(Z_r\) has no pole at its generic divisor, by birational
comparison. Its scalar coefficient \(u\), pulled back by \(t=x^e\),
therefore extends meromorphically, with
\[
 e\,\ord_z(u)+m\bigl(rh-(r-1)(e-1)\bigr)\geq0.
\]
In particular,
\begin{equation}\label{pos:fiber-power-pole-bound}
 \ord_z(u)\geq-\frac{mrh}{e}.
\end{equation}
There is no pole allowance at a nonexceptional base curve, where
\(h=0\). The exceptional allowances are exactly
\(mr c_P\), with the fixed \(c_P\) in
\eqref{pos:exceptional-pole-allowance}.

The nonzero meromorphic coefficient in
\eqref{pos:fiber-power-coefficient-line}, with these pole bounds,
implies
\[
 \deg K_C+r\left(A-K_S+\sum_Pc_PP\right)\cdot C\geq0.
\]
Divide by \(r\) and let \(r\to\infty\), keeping \(C\) fixed.
Since \(C\) is a positive multiple of \(H'\), this proves
\eqref{pos:surface-intersection}.
\end{proof}

Take the surface Zariski decomposition \cite{BHPV}
\begin{equation}\label{pos:zariski-decomposition}
 A=P_0+N_0,
\end{equation}
where \(P_0\) is nef, \(N_0\geq0\) has negative-definite support
if nonzero, and \(P_0\) is orthogonal to every component of \(N_0\).
The negative-part linear equations give rational coefficients, so
\(P_0=A-N_0\) remains an actual rational line. For a
\(\mu\)-exceptional curve \(P\), \(A\cdot P=0\). If \(P\) belongs
to \(\Supp N_0\), orthogonality gives \(P_0\cdot P=0\); otherwise
both \(P_0\cdot P\) and \(N_0\cdot P\) are nonnegative, and their
sum is zero. Hence every such \(P\) is orthogonal to \(P_0\).

If \(P_0\equiv0\), the equality \(q(S)=0\) makes a divisible
multiple of \(P_0\) torsion, so \(P_0\) is zero in
\(\Pic(S)\otimes\Q\). If \(P_0^2>0\), the nef line \(P_0\) is
big. Either case supplies sections of a positive multiple of \(A\).
In the remaining case, \(P_0\not\equiv0\) and \(P_0^2=0\).
Test \eqref{pos:surface-intersection} on ample classes approaching
\(P_0\). The exceptional terms and \(A\cdot P_0\) vanish, giving
\(K_S\cdot P_0\leq0\). For large divisible \(\ell\),
\begin{equation}\label{pos:surface-riemann-roch}
 \chi(S,\ell P_0)
 =\chi(\OO_S)-\frac{\ell}{2}K_S\cdot P_0>0,
\end{equation}
since \(\chi(\OO_S)=1+h^{2,0}(S)>0\). Serre duality gives
\(H^2(S,\ell P_0)=0\) for large \(\ell\): the divisor
\(K_S-\ell P_0\) has negative degree against a fixed ample divisor.
Thus \(h^0(S,\ell P_0)>0\); multiplication by the section of
\(\ell N_0\) proves \(\kappa(S,A)\geq0\).

\subsection{A threefold base and genus-one orders}

Assume \(d=3\). The general fibers of \(g:V\to T\) are smooth
connected genus-one curves. Indeed, the singular locus of normal
\(V\) has dimension at most two, so misses a general fiber; generic
smoothness then applies, and \eqref{pos:canonical-pullback} makes its
canonical line torsion. Also all components of \(R\) in
\eqref{pos:resolved-pullback} are vertical over \(S\): their images
on \(V\) have dimension at most two.

Choose one integer \(m>0\), divisible by \(12\), that clears the
Cartier data and the actual isomorphism in
\eqref{pos:canonical-pullback}. It clears
\eqref{pos:resolved-pullback} on every smooth model subsequently
used. In fact, the pullback of a local frame of the Cartier line
\(\OO_V(mK_V)\), regarded on the regular locus as an
\(m\)-pluriform, has integral orders on every resolution. Those
orders are precisely \(m\operatorname{coeff}_E R\). There is thus
no need to choose a new index after a base blowup.

\subsubsection{The actual Hodge line and its growth}

Shrink to a dense open \(U\subset T_{\mathrm{reg}}\), also viewed
on \(S\), on which the family is smooth and agrees with its resolution.
Let \(\mathcal H\) be its line of holomorphic fiberwise one-forms.
It is a holomorphic line by Grauert base change; periods against local
integral cycles are holomorphic. Give it the Hodge norm, so that
\(\|\alpha\|^2\) is, up to a fixed convention, the fiber integral
of \(\ii\alpha\wedge\overline\alpha\).

Evaluation \(f^*\mathcal H\to\omega_{M/U}\) is an isomorphism:
a nonzero holomorphic one-form on a smooth genus-one curve has no zero.
Its \(m\)-th power, together with the fixed actual pullback identity,
gives an isomorphism of pullbacks of lines on \(U\). Applying
\(f_*\) and \(f_*\OO_M=\OO_U\) yields the actual identification
\begin{equation}\label{pos:actual-hodge-line}
 \mathcal H^{\otimes m}\simeq\mathcal L_m|_U,
 \qquad\mathcal L_m=\OO_S\bigl(m(A-K_S)\bigr).
\end{equation}
It is compatible on all birational base models over their common
open. No section of the genus-one fibration is needed: its first
integral cohomology, periods, and invariant differentials are defined
without an origin. In particular, \eqref{pos:actual-hodge-line}
leaves no unspecified flat line factor.

Let \(E_4\) and \(\Delta_{\mathrm{mod}}\) be the classical modular
forms of weights four and twelve for \(\mathrm{SL}_2(\Z)\)
\cite[Chapter~VII]{SerreArithmetic}. A local symplectic period basis
gives a parameter \(\tau\) in the upper half-plane and a normalized
Hodge frame with periods \(1,\tau\). The weight transformation laws
compensate the frame transformation, so these forms define sections
of \(\mathcal H^4\) and \(\mathcal H^{12}\). Define sections of
\(\mathcal L_m|_U\) by
\begin{equation}\label{pos:modular-sections}
 e_1=E_4^{m/4},\qquad e_2=\Delta_{\mathrm{mod}}^{m/12},
\end{equation}
where the tensor powers of the normalized frame are understood.

Near a general point of a complementary prime divisor on any smooth
base model, let \(t=0\) be its equation. There are positive constants
\(c,C,b\), locally uniform in the remaining coordinates, such that
\begin{equation}\label{pos:modular-growth}
 c\leq\bigl(\|e_1\|+\|e_2\|\bigr)^{2/m}
 \leq C(1+|\log|t||)^b\qquad(t\ne0).
\end{equation}
Here is the full growth argument. In the standard fundamental domain
the normalized frame has squared norm proportional to
\(\operatorname{Im}\tau\). Both scalar modular forms are bounded
there, \(\Delta_{\mathrm{mod}}\) has no zero in the upper half-plane,
and \(E_4\to1\) at the cusp. On the compact part their joint norm has
a positive minimum; at the cusp the \(E_4\) term supplies the same
positive lower bound. This also covers the case where \(e_1\) is
identically zero on the given family.

For the upper bound, restrict to punctured disks transverse to the
prime, with the other parameters in a small compact set and with
\(t\ne0\) inside \(U\). Lift the period map to universal covers.
Schwarz--Pick makes it distance-decreasing for the hyperbolic metrics.
The radial hyperbolic distance in a punctured disk, from a fixed
radius to \(|t|\), is
\(O(1)+O(\log|\log|t||)\). Starting period representatives at that
radius can be chosen in a bounded part of the fundamental domain,
uniformly in angle and the other parameters, by compactness inside
the smooth locus. Since \(\log\operatorname{Im}\tau\) changes by
at most hyperbolic distance, both \(\operatorname{Im}\tau\) and its
inverse are bounded by powers of \(1+|\log|t||\) along the lifted
radial paths. Returning to the fundamental domain preserves such a
bound, because for \(\gamma\in\mathrm{SL}_2(\Z)\),
\[
 \operatorname{Im}(\gamma\tau)
 \leq\max\{\operatorname{Im}\tau,(\operatorname{Im}\tau)^{-1}\}.
\]
The bounded scalar modular forms and the Hodge-frame norm now prove
the upper inequality in \eqref{pos:modular-growth}.

\subsubsection{The integration threshold and meromorphic extension}

Fix a prime \(P\subset S\), a local parameter \(t\) at its general
point, and a frame \(s\) of \(\mathcal L_m\). On a simultaneous log
resolution of \(R\) and \(f^*P\), recompute \(R\) by the canonical
comparison in \eqref{pos:resolved-pullback}. Define
\begin{equation}\label{pos:threshold-definition}
 t_P=\min_{E\mapsto P}
 \frac{1+\operatorname{coeff}_E R}{\ord_E(f^*P)}.
\end{equation}
The minimum is taken over components above the general point of
\(P\). It is positive because \(R\geq0\).

\begin{lemma}[The divisorial order identity]\label{pos:order-identity}
The sections \(e_1,e_2\) extend meromorphically to \(S\). If
\[
 \ell_P=\min_i\ord_P(e_i/s),
\]
omitting an identically zero section, then
\begin{equation}\label{pos:modular-order}
 \frac{\ell_P}{m}=1-t_P.
\end{equation}
The same identity holds on every subsequent smooth base model, using
its actual line \(\OO(m(A-K_S))\) and the fixed integer \(m\).
\end{lemma}

\begin{proof}
First identify the weighted integrability threshold of the frame:
\begin{equation}\label{pos:integrability-threshold}
 t_P=\sup\left\{u\in\R:
 |t|^{-2u}\|s\|^{2/m}\text{ is locally integrable near general }P
 \right\}.
\end{equation}
Multiply \(s\) by a local \(m\)-canonical base frame. Under
\eqref{pos:resolved-pullback}, the resulting total-space
pluriform has divisor \(mR\). On a smooth good fiber it is the
\(m\)-th tensor of a one-form, up to a scalar, so integrating its
\(2/m\)-density along the fiber gives \(\|s\|^{2/m}\) times the
smooth base coordinate density. Fubini and change of variables
therefore identify the weighted base integral with the total-space
integral. In simple normal crossing coordinates the latter has, at
a component \(E\), a factor
\[
 |x_E|^{2(h_E-u e_E)},\qquad
 h_E=\operatorname{coeff}_E R,\quad e_E=\ord_E(f^*P).
\]
Such a factor is integrable exactly when \(h_E-u e_E>-1\).
Choose the point of \(P\) generally to exclude components not
dominating it, and use properness for a finite covering of its inverse
image. The conditions are exactly those in
\eqref{pos:threshold-definition}, proving
\eqref{pos:integrability-threshold}.

We next establish a polynomial lower bound on the Hodge norm of
\(s\), without assuming meromorphic extension of the modular
coefficients. Choose one component \(E\) above general \(P\), of
multiplicity \(e\) and coefficient \(h\) in \(R\). At a general
point of \(E\), the map has local coordinates
\[
 (t,z_1,z_2)=(x^e,z_1,z_2)
\]
on the base, with one further fiber coordinate \(y\) on the source.
The map \(E\to P\) is submersive there, units have been absorbed
into \(x\), and no other divisorial support meets the chart. The
total pluriform associated with \(s\) is a unit times
\(x^{mh}(dx\wedge dz_1\wedge dz_2\wedge dy)^m\).
After division by the base pluriform, its relative coefficient is
a unit times \(x^{m(h-e+1)}\). Integrating on a fixed smaller disk
in the \(y\)-coordinate gives
\begin{equation}\label{pos:polynomial-lower-bound}
 \|s\|^{2/m}\geq c\,|t|^{2(h-e+1)/e}.
\end{equation}
These charts are available above every point of \(P\) outside a
proper analytic subset. Indeed, the non-submersive locus and
intersections with the removed supports are proper subsets of \(E\);
those that dominate \(P\) have strictly smaller generic fiber
dimension. Properness and the fiber-dimension theorem show that a
general fiber of \(E\to P\) is not exhausted by them. We also omit
the singular and intersection loci of complementary base divisors.

On the punctured chart, write \(e_i=a_i s\). Combining
\eqref{pos:modular-growth} with
\eqref{pos:polynomial-lower-bound} gives an upper bound for
\(|a_i|\) by a fixed power of \(|t|^{-1}\), after absorbing a
logarithmic power into an arbitrarily small additional power. Thus
\(a_i\) has at most a pole across general \(P\). The integer pole
allowance can be fixed for that prime because the exponent in
\eqref{pos:polynomial-lower-bound} depends on the one fixed component
\(E\). There are finitely many complementary prime divisors. Twist
\(\mathcal L_m\) by their bounded pole allowances; both sections
are then holomorphic away from an analytic set of codimension at
least two. Hartogs extension on the smooth \(S\) gives global
meromorphic sections of the original \(\mathcal L_m\).

At a general point of \(P\), the two meromorphic scalar coefficients
satisfy
\[
 |a_1|+|a_2|\asymp |t|^{\ell_P}.
\]
Equation~\eqref{pos:modular-growth} therefore gives
\begin{equation}\label{pos:frame-order-growth}
 c|t|^{-2\ell_P/m}
 \leq\|s\|^{2/m}
 \leq C|t|^{-2\ell_P/m}(1+|\log|t||)^b.
\end{equation}
The supremum of the exponents \(u\) for which the weighted expression
is integrable is consequently \(1-\ell_P/m\): it is integrable for
strictly smaller \(u\), and the lower bound makes it nonintegrable
for strictly larger \(u\). Its possible behavior at the endpoint
does not change the supremum. Comparison with
\eqref{pos:integrability-threshold} proves
\eqref{pos:modular-order}. At a good divisor the identity is immediate.

On a further smooth base modification the same argument uses its
canonical line and a simultaneous resolution of the total space.
The integer \(m\) still clears the actual identity, as explained
above. The sections agree on the common open by
\eqref{pos:actual-hodge-line}, and hence over the meromorphic field.
Thus the calculation applies with the same \(m\) on every model
needed below.
\end{proof}

For a prime \(P\) that is not exceptional over \(T\), one also has
\begin{equation}\label{pos:nonexceptional-threshold}
 t_P\leq1.
\end{equation}
Indeed, as in the surface argument, a component of the inverse image
of its image divisor on \(T\) is a divisor on \(V\) dominating it.
Its strict transform has coefficient zero in \(R\) and positive
integral pullback multiplicity. It is one of the components tested in
\eqref{pos:threshold-definition}. No such upper bound is required
for a base-exceptional prime.

\subsubsection{A klt adjoint on the original base}

Resolve the meromorphic pencil \([e_1:e_2]\) by projective blowups
of \(S\), and then resolve the divisorial supports, including the
exceptional locus over \(T\). A constant pencil is allowed. If
\(\rho:S'\to S\) is such a smooth modification, its actual line is
\begin{equation}\label{pos:base-canonical-transform}
 \mathcal L'_m
 =\rho^*\mathcal L_m\otimes\OO_{S'}(-mK_{S'/S}).
\end{equation}
Consequently both section divisors transform by their pullbacks minus
the same relative canonical Jacobian term. Resolving the pencil and
subtracting its full signed fixed divisor therefore leaves a free
moving system; subsequent blowups pull back that free system and
preserve its freeness.

Rename this smooth model \(S\). Its fixed divisor is
\(\sum_P\ell_PP\), with simple normal crossing support. A general
complex linear combination \(e\) of \(e_1,e_2\) has
\begin{equation}\label{pos:crepant-boundary-upstairs}
 \frac1m\divisor_S(e)
 =\sum_P(1-t_P)P+\frac1mQ=:B_S,
\end{equation}
where \(Q\) is a general member of the free moving system, possibly
zero. Bertini makes it smooth, reduced, and transverse to all the
chosen strata, with no fixed component. Every coefficient of this
simple normal crossing divisor is strictly less than one: the fixed
coefficients are \(1-t_P<1\), and the moving coefficient is
\(1/m<1\). The fixed coefficients over exceptional primes may be
negative.

Regard \(e\) downstairs as a rational section of the rank-one
reflexive difference \(mA_T-mK_T\), and define
\begin{equation}\label{pos:base-boundary}
 \Xi=\frac1m\divisor_T(e).
\end{equation}
Equation~\eqref{pos:nonexceptional-threshold} shows that
\(\Xi\geq0\). There is an actual rational-linear identity
\begin{equation}\label{pos:base-adjoint}
 K_T+\Xi\sim_{\Q}A_T.
\end{equation}
To verify both this identity and crepancy, choose over the rational
function field a frame \(a\) of \(mA_T\) and a rational canonical
form \(\omega\), and write \(e=u\,a/\omega^m\). With canonical
divisor chosen by \(\omega\),
\[
 m(K_T+\Xi)=\divisor(u)+\divisor(a).
\]
The right side is Cartier. Pull it back to \(S\) and subtract
\(m\divisor_S(\omega)\); the result is precisely
\(\divisor_S(e)\), with the canonical Jacobian term in
\eqref{pos:base-canonical-transform}. Thus
\begin{equation}\label{pos:actual-crepant-identity}
 K_S+B_S=\mu^*(K_T+\Xi)
\end{equation}
for compatible rational divisor representatives. This proves that
\((T,\Xi)\) is klt: \(B_S\) is its simple normal crossing crepant
boundary with every coefficient less than one. Only the sum
\(K_T+\Xi\) has been proved \(\Q\)-Cartier; a separate
\(\Q\)-Gorenstein hypothesis on \(T\) was not used.

Finally put \(\Gamma_S=(B_S)_+\). It is an effective rational klt
simple normal crossing boundary, and
\[
 K_S+\Gamma_S\sim_{\Q}A+E_S,
 \qquad E_S\geq0\text{ and }E_S\text{ is }\mu\text{-exceptional}.
\]
The exceptionality follows from \(\Xi\geq0\): negative coefficients
of the crepant boundary occur only over \(T\)'s codimension-two
locus. This adjoint is numerically pseudo-effective on the smooth
projective threefold \(S\). Projective klt threefold nonvanishing,
in the scope recalled in Section~\ref{red:section}, gives a section
of a positive multiple of \(K_S+\Gamma_S\). Push it to \(T\).
The exceptional pole allowance disappears in codimension one, and
normality extends the section of the actual line \(mA_T\) across
codimension two, after taking a common multiple. Hence
\(\kappa(T,A_T)\geq0\), as required.

\begin{proof}[Proof of Proposition~\ref{pos:nonvanishing}]
Proposition~\ref{pos:pullback-model} gives
\eqref{pos:resolved-pullback} with \(A\) pseudo-effective and
\(q(S)=0\). The curve, surface, and threefold arguments above each
give \(\kappa(S,A)\geq0\) in the corresponding dimension. Pull a
divisible section back to \(M\) and multiply by the section of
\(mR\) in \eqref{pos:resolved-pullback}. This gives a nonzero
canonical plurisection on \(M\), and smooth birational invariance
gives \(\kappa(W,K_W)\geq0\).
\end{proof}
